\documentclass[11pt,a4paper]{article}
\usepackage[T1]{fontenc}
\usepackage{lmodern,amsmath,amssymb,amsthm,mathtools,microtype}
\usepackage[margin=27mm]{geometry}
\usepackage[colorlinks=true,linkcolor=blue,citecolor=blue,urlcolor=blue]{hyperref}
\setlength{\emergencystretch}{2em}
\newtheorem{theorem}{Theorem}[section]
\newtheorem{proposition}[theorem]{Proposition}
\newtheorem{corollary}[theorem]{Corollary}
\theoremstyle{remark}\newtheorem{remark}[theorem]{Remark}
\newcommand{\Q}{\mathbb Q}\newcommand{\PP}{\mathbb P}
\DeclareMathOperator{\N}{N}
\title{A rational quadric model for a varying-input pole surface}
\author{}\date{16 November 2026}
\begin{document}\maketitle
\begin{abstract}
Fix a positive rational weight and a nondegenerate rational point on the
first pole conic. When the second pole input varies in its prescribed
square class, its surface is birational over the rationals to a smooth
quadric with an explicit rational point. We give mutually inverse maps,
a rational parametrization and a rational curve with nonzero ordinates
along which the input tends to zero through positive values. Products
with a common first-conic parameter are rational as well. The input scale
is free in these statements; imposing a separately prescribed elliptic
detector expression for that scale remains an additional problem.
\end{abstract}

\section{The surface and its rational point}
Fix $D\in\Q_{>0}$ and $\tau\in\Q$ satisfying
\begin{equation}\label{eq:hypotheses}
                       \tau\ne0,\qquad \tau^2\ne D.
\end{equation}
Put
\begin{equation}\label{eq:ep}
 e=\frac{\tau^2-D}{\tau^2+3D},\qquad
 p=-\frac{4D\tau}{\tau^2+3D}.
\end{equation}
Then
\begin{equation}\label{eq:conic}
                   p^2=D(1-e)(1+3e),
\end{equation}
and $-1/3<e<1$, $e\ne0$, $p\ne0$. Conversely, every rational
pair satisfying~\eqref{eq:conic} with $D>0$, $e\ne0$ and $p\ne0$
has~\eqref{eq:ep} with $\tau=-p/(1-e)$.
The excluded first-conic cases are explicit. If $e=0$, then $D=p^2$
and the surface below reduces to $Q^2=Du^4$. If $p=0$, then
$e=1$ or $e=-1/3$; the designated first ordinate vanishes. Neither
case is used in the nonzero-ordinate statements.

Write $w=D/68$. Let $V\in\Q^\times$ be free and set
\begin{equation}\label{eq:input}
                 b=\frac{4w}{V^2}=\frac D{17V^2},
                 \qquad C=17(17-b)=289-\frac D{V^2}.
\end{equation}
The normalized second pole equation is
\begin{equation}\label{eq:normalized}
                         q^2=D(s^2-eC)(s^2+3eC).
\end{equation}
With $u=sV$ and $Q=qV^2$, this becomes the affine surface
\begin{equation}\label{eq:surface}
 \mathcal S_{D,e}:\quad
 Q^2=D\bigl(u^2-e(289V^2-D)\bigr)
          \bigl(u^2+3e(289V^2-D)\bigr).
\end{equation}
Its weighted projective closure in $\PP(1,1,1,2)$ is
\begin{equation}\label{eq:projective}
 Q^2=D(u^2-289eV^2+eDT^2)(u^2+867eV^2-3eDT^2).
\end{equation}
The two branch conics are smooth and distinct over an algebraic closure,
so their product is not a square; the surface is geometrically integral.
The rational point
\[
                     (u:V:T:Q)=(17:1:0:289p)
\]
lies off the branch divisor because $p\ne0$.

\section{Exact birational maps}
Consider the affine quadric
\begin{equation}\label{eq:quadric}
 \mathcal Q_{D,e}:\quad x^2-Dy^2=e(\sigma^2+3D).
\end{equation}
Its projective closure is smooth: its diagonal quadratic form has
nonzero determinant $-3e^2D^2$. It contains the rational point
\begin{equation}\label{eq:point}
                         P_0=(\tau,1,\tau).
\end{equation}

\begin{theorem}\label{thm:maps}
Put
\[
 R=289V^2-D,\qquad A=u^2-eR,
 \qquad N=x^2-Dy^2.
\]
There is an isomorphism between the dense open sets
\begin{equation}\label{eq:opens}
 \mathcal S^\circ:\ uRA\ne0,
 \qquad
 \mathcal Q^\circ:\
      (\sigma^2+3D)(\sigma^2-D)(x-\sigma y)\ne0.
\end{equation}
The map from $\mathcal Q^\circ$ to $\mathcal S^\circ$ is
\begin{equation}\label{eq:forward}
 \begin{gathered}
 \alpha=\frac{\sigma x-Dy}{N},\qquad
 \beta=\frac{x-\sigma y}{N},\qquad
 r=-\frac{4D\sigma}{\sigma^2+3D},\\
 u=\frac1\beta,\qquad
 V=\frac\alpha{17\beta},\qquad
 Q=\frac r{\beta^2}.
 \end{gathered}
\end{equation}
The inverse is
\begin{equation}\label{eq:reverse}
 \sigma=-\frac QA,\qquad
 x=\frac{u(17\sigma V-D)}R,\qquad
 y=\frac{u(17V-\sigma)}R.
\end{equation}
The additional conditions $V\ne0$ and $Q\ne0$ correspond respectively
to $\sigma x-Dy\ne0$ and $\sigma\ne0$.
\end{theorem}

\begin{proof}
Use the quadratic algebra $\Q[\xi]/(\xi^2-D)$, including the split
case when $D$ is a square. Its norm is
$\N(a+b\xi)=a^2-Db^2$, and an element is invertible exactly when
its norm is nonzero.

On $\mathcal S^\circ$, define
\[
 z=\frac{eR}{u^2},\qquad r=\frac Q{u^2},\qquad
 g=\frac{17V+\xi}{u}.
\]
Equation~\eqref{eq:surface} gives
\[
 r^2=D(1-z)(1+3z),\qquad \N(g)=z/e.
\]
Here $z\ne0$ and $z\ne1$. Parametrizing this conic by
$\sigma=-r/(1-z)$ gives
\begin{equation}\label{eq:conicparameter}
 z=\frac{\sigma^2-D}{\sigma^2+3D},\qquad
 r=-\frac{4D\sigma}{\sigma^2+3D}.
\end{equation}
Consequently, if
\[
                         x+y\xi=\frac{\sigma+\xi}{g},
\]
then
\[
 x^2-Dy^2=\frac{\sigma^2-D}{z/e}
                        =e(\sigma^2+3D).
\]
Expanding the quotient gives~\eqref{eq:reverse}.

Conversely, on $\mathcal Q^\circ$ the denominator
$N=e(\sigma^2+3D)$ is nonzero. Set
\[
               g=\frac{\sigma+\xi}{x+y\xi}
                 =\alpha+\beta\xi.
\]
Its norm is $(\sigma^2-D)/N$. Since $\beta\ne0$, equations
\eqref{eq:forward} give $g=(17V+\xi)/u$. Moreover
\[
 \frac{e(289V^2-D)}{u^2}
       =e\N(g)=\frac{\sigma^2-D}{\sigma^2+3D}.
\]
The conic identity~\eqref{eq:conicparameter} now proves
\eqref{eq:surface}.

For the exact correspondence of open sets, the forward identities are
\[
 R=\frac{\sigma^2-D}{N\beta^2},\qquad
 A=\frac{4D}{(\sigma^2+3D)\beta^2}.
\]
In the reverse direction one has
\[
 \sigma^2+3D=\frac{4Du^2}{A},\qquad
 \sigma^2-D=\frac{4DeR}{A},\qquad
 x-\sigma y=\frac{4Deu}{A}.
\]
Thus every required denominator is nonzero. The two quotient
constructions recover $g$ and $x+y\xi$ mutually, and the conic
parametrization recovers $z,r$ mutually. This proves both inverse
identities. Finally~\eqref{eq:forward} gives the asserted conditions
for nonzero $V$ and $Q$.
\end{proof}

\section{Rational parametrization}
The rational point~\eqref{eq:point} supplies an explicit parametrization
of the smooth quadric. For a direction $(a:b:c)\in\PP^2$, put
\[
 H=a^2-Db^2-ec^2,\qquad
 L=\tau a-Db-e\tau c.
\]
On $HL\ne0$, define
\begin{equation}\label{eq:projection}
 k=-\frac{2L}{H},\qquad
             (x,y,\sigma)=(\tau,1,\tau)+k(a,b,c).
\end{equation}
This is invariant under rescaling the direction. Substitution into
the quadric equation leaves $2kL+k^2H=0$. Conversely, for a quadric
point other than $P_0$, the direction is
\[
                         (x-\tau:y-1:\sigma-\tau).
\]
On the dense open set where its quadratic value $H$ is nonzero,
the quadric equation gives $H+2L=0$, so~\eqref{eq:projection}
returns that point. Hence the projection is birational.

\begin{corollary}\label{cor:rational}
The surface $\mathcal S_{D,e}$ and its projective closure are rational
over $\Q$. Equations~\eqref{eq:projection} and~\eqref{eq:forward}
give an explicit dominant rational parametrization by two parameters.
\end{corollary}

The point $P_0$ itself corresponds to $\beta=0$ and therefore to
the point at $T=0$ in~\eqref{eq:projective}. The parametrization uses
the open loci specified above to obtain finite affine points; it does
not treat that boundary point as an affine witness.

\section{Nonzero ordinates and small positive inputs}
The following curve gives a direct affine construction, independently
of choices of projection directions.

\begin{proposition}\label{prop:curve}
For $h\in\Q^\times$ with $Dh^2\ne1$, set
\begin{equation}\label{eq:curve}
 u=\frac{1-Dh^2}{2h},\qquad
 V=\frac{1+Dh^2}{34h},\qquad Q=p u^2.
\end{equation}
These points belong to $\mathcal S^\circ$, with $u,V,Q\ne0$.
Their native pole coordinates and input are
\begin{equation}\label{eq:nativecurve}
 \begin{gathered}
 b=\frac{68Dh^2}{(1+Dh^2)^2},\qquad
 s=\frac{17(1-Dh^2)}{1+Dh^2},\\
 C=s^2,\qquad q=p s^2.
 \end{gathered}
\end{equation}
In particular $0<b<17$, $s,q\ne0$, and $b$ has the rational square
class of $w$. As $h\to0$ through nonzero rational values,
\[
                 b\longrightarrow0,\qquad
                 s\longrightarrow17,\qquad q\longrightarrow289p.
\]
Thus there are infinitely many finite affine witnesses with $0<b<1$
and all displayed pole ordinates nonzero.
\end{proposition}

\begin{proof}
The identity $289V^2-D=u^2$ follows from~\eqref{eq:curve}.
Thus the right side of~\eqref{eq:surface} is
$D(1-e)(1+3e)u^4=p^2u^4=Q^2$. Also
$R=u^2\ne0$ and $A=(1-e)u^2\ne0$, proving the open conditions.
Since $D>0$ and $p\ne0$, none of $V,Q$ vanishes.

Substitution gives~\eqref{eq:nativecurve} and the exact identity
\[
             17-b=17\frac{(1-Dh^2)^2}{(1+Dh^2)^2}>0.
\]
The square-class assertion is immediate from $b=4w/V^2$.
The limits follow from the displayed rational functions. For example,
$0<b\le68Dh^2$, so every sufficiently small nonzero rational $h$
gives $b<1$. There are infinitely many such values, and the nonconstant
function $V=(1+Dh^2)/(34h)$ has finite fibers.
\end{proof}

The proposition verifies the real size condition and the rational square
class explicitly. Additional finite-prime preparation is a separate
condition on the actual input; it is not inferred from real smallness.

\section{A common first-conic parameter and the remaining lifting problem}
Fix positive rational weights $w_1,\ldots,w_m$, put $D_i=68w_i$,
and choose a common rational $\tau\ne0$ with $\tau^2\ne D_i$ for all
$i$. Define $e_i,p_i$ by~\eqref{eq:ep} with $D=D_i$.

\begin{corollary}\label{cor:product}
With each $V_i$ free, the product
\[
                         \prod_{i=1}^m\mathcal S_{D_i,e_i}
\]
is rational over $\Q$, of dimension $2m$. Its simultaneous native
pole equations admit nonzero rational witnesses with every
$0<b_i<1$. If the common $\tau$ also varies, restrict the base by
\[
                \tau\prod_{i=1}^m(\tau^2-D_i)(\tau^2+3D_i)\ne0.
\]
The corresponding total space is rational of dimension $2m+1$.
\end{corollary}
\begin{proof}
Each factor has the explicit birational parametrization by $\PP^2$,
and the parameters for different factors are independent. Applying
Proposition~\ref{prop:curve} with independently small rational $h_i$
gives the simultaneous real-size assertion with the same fixed $\tau$.
The formulas are rational functions of $\tau$ as well. Over $\Q(\tau)$
they give the same birational maps, so the total space is birational
to an open subset of $\mathbb A^1\times(\PP^2)^m$.
\end{proof}

The freedom to choose $V_i$ is part of these statements. If an elliptic
detector prescribes $V_i$ as a rational function of other witness
coordinates, one must also solve those identities. Rationality of the
free-input surface does not establish that pullback, a solution with
arbitrarily prescribed $V_i$, or a reduction in the number of variables
of a complete defining polynomial. Keeping $b$ fixed instead gives
the different surface studied in~\cite{Fixed,Sections}.

\section{Exact certificate}
The package \nolinkurl{computation/pole_parameters/varying_input_quadric_v44}
contains a standalone exact checker and its complete output. It verifies
the conic identities, the maps in both directions, the exact matching
of open loci, the projection parametrization, and the native coordinates
of the positive curve. Polynomial remainders are taken modulo the
appropriate surface or quadric equation. An independent rational
instance and two sign-error controls supplement these symbolic identities.
The certificate addresses the stated algebraic maps; it does not check
any prescribed detector pullback.

\begingroup\raggedright
\begin{thebibliography}{9}
\bibitem{Fixed}\emph{K3 surfaces in the fixed-input pole problem},
project manuscript, 16 November 2026,
\nolinkurl{negative_results/papers/pole_parameters/fixed_input_k3_surfaces_20261116_v42.tex}.
\bibitem{Sections}\emph{Rational sections of fixed-input pole fibrations},
project manuscript, 16 November 2026,
\nolinkurl{negative_results/papers/pole_parameters/rational_sections_20261116_v42.tex}.
\end{thebibliography}
\endgroup
\end{document}
